% Deux arbres de Merkle de 8 feuilles qui ne diffèrent que par la feuille 5 :
% la comparaison suit le seul chemin dont les hachages diffèrent.
\documentclass[border=6pt]{standalone}
\usepackage{sd}
\begin{document}
\begin{tikzpicture}[x=1mm, y=1mm, note/.append style={font=\sffamily\small},
    nd/.style={draw=#1, fill=#1Bg, line width=0.7pt, minimum width=13mm, minimum height=7mm,
      inner sep=0pt, font=\ttfamily\scriptsize, rounded corners=2pt},
    nd/.default=Dslate,
    lab/.style={font=\sffamily\small, text=#1, align=center},
    lab/.default=Dink]
  \newcommand{\arbre}[3]{% #1 : décalage x, #2 : nom, #3 : hachages différents (1) ou non (0)
    \begin{scope}[xshift=#1 mm]
      \node[lab, font=\sffamily\small\bfseries] at (35,12) {#2};
      % racine
      \ifnum#3=1 \node[nd=Dred] (r) at (35,0) {c41a…}; \else \node[nd] (r) at (35,0) {9e07…};\fi
      % niveau 1
      \node[nd] (n0) at (17.5,-14) {3b2f…};
      \ifnum#3=1 \node[nd=Dred] (n1) at (52.5,-14) {d802…}; \else \node[nd] (n1) at (52.5,-14) {71ac…};\fi
      % niveau 2
      \foreach \i/\x/\h in {0/8.75/a4e1, 1/26.25/5c90, 3/61.25/e3b7} \node[nd] (m\i) at (\x,-28) {\h…};
      \ifnum#3=1 \node[nd=Dred] (m2) at (43.75,-28) {0f6d…}; \else \node[nd] (m2) at (43.75,-28) {82c4…};\fi
      % feuilles
      \foreach \i in {0,...,7} {
        \pgfmathsetmacro{\x}{4.375+\i*8.75}
        \ifnum\i=5 \ifnum#3=1 \node[nd=Dred, minimum width=8mm] (f\i) at (\x,-42) {\i}; \else \node[nd, minimum width=8mm] (f\i) at (\x,-42) {\i};\fi
        \else \node[nd, minimum width=8mm] (f\i) at (\x,-42) {\i};\fi
      }
      \draw[Dslate] (r) -- (n0) (r) -- (n1);
      \draw[Dslate] (n0) -- (m0) (n0) -- (m1) (n1) -- (m2) (n1) -- (m3);
      \foreach \p/\a/\b in {m0/0/1, m1/2/3, m2/4/5, m3/6/7} \draw[Dslate] (\p) -- (f\a) (\p) -- (f\b);
    \end{scope}}
  \arbre{0}{réplique A}{0}
  \arbre{85}{réplique B (a manqué une écriture)}{1}
  % étapes
  \node[lab=Dred, anchor=west, align=left] at (160,0) {1 : racines différentes};
  \node[lab=Dred, anchor=west, align=left] at (160,-14) {2 : seul le fils droit diffère};
  \node[lab=Dred, anchor=west, align=left] at (160,-28) {3 : seul le nœud 2 diffère};
  \node[lab=Dred, anchor=west, align=left] at (160,-42) {4 : feuille 5 à comparer};
  \node[lab=Dslate, anchor=north west, align=left] at (0,-52) {une feuille = le XOR des hachages (clé, valeur) des clés de son intervalle ;
    un nœud = le hachage de ses deux enfants ;\\un aller-retour par niveau, et seuls les sous-arbres dont les hachages diffèrent sont parcourus};
\end{tikzpicture}
\end{document}
